\scanpage{25}\mspage{24}
Finally, we shall work with the topology defined first (basis formed with the ``intervals''), and we shall nevertheless call, by abuse of language, the $[a,b]$, $\mathcal L_{\leq a}$, $\mathcal L_{\geq b}$ (as well as $\mathcal L$) ``closed'' \unsure{ordered|ordinary} (or \unsure{biordered|well-given|well let us say|binominate}) \uline{sub-segment pieces}, [NB we assume them not reduced to a point, i.e.
\[
\deleted{\unsure{[|i}} a<b\quad \text{resp.}\quad a\ne\mathrm{or}(\mathcal L)\quad \text{resp.}\quad \deleted{\unsure{a}}b\ne\mathrm{ex}(\mathcal L)\;]
\]

A \uline{broad \inserted{closed} sub-segment piece} is a \deleted{sub-segment piece} subset which is either a sub-segment piece, or \unsure{reduced|reduced} to a point.

For such a subset, $T$, we
% Widely spaced strokes cross the whole lower block, from pose to the last example.
% The close zigzags and local horizontal cancellations show an earlier level of correction.
\deleted{set}
\[
\begin{aligned}
\deleted{\partial T=}&\deleted{\{t\in T\mid t\in \overline{\mathcal L\setminus T}\}}\quad\text{\deleted{boundary}}\\
 &\text{\deleted{\unsure{boundary in the topology|boundary in the topology}}}\\
 &\text{\deleted{absolute of $T$ $=$ set}}\\
 &\text{\deleted{canonical}}\\
 \deleted{={}\unsure{T\cap\mathcal L|T\cap\overline{\mathcal L}}}&\quad\text{\deleted{closed of its endpoints}}\\
 &\deleted{(\mathrm{NB}\ \mathrm{card}\,\partial T\in\{0,1,2\})}
\end{aligned}
\]

\deleted{We see that for a closed segment piece (\uline{NB} not reduced to a point) we have}

\[
\begin{aligned}
 \deleted{\partial[a,b]=\{a,b\}}&\qquad \deleted{\partial \mathcal L_{\leq a}=\{a\},\quad \partial \mathcal L_{\geq a}=\{a\}}\\
 &\deleted{\partial\mathcal L=\emptyset}
\end{aligned}
\]

% faint vertical paper watermark
% the paper's watermark (EXTRA STRONG), not his text
